% !TEX program = xelatex
% ============================================================
%  سجل القرارات — Decision Log / ADR Template
%  نموذج لتوثيق قرارات المشروع والبدائل المطروحة وسياق كل قرار
%  Compile with: xelatex main.tex (run twice)
% ============================================================
%  Features:
%    - Decision table: ID, Date, Decision, Context, Alternatives,
%      Decision maker, Status
%    - 5 sample decision rows
%    - Detailed decision records section
% ============================================================

\documentclass[11pt,a4paper]{article}

\usepackage[a4paper,margin=2.5cm,top=2.5cm,bottom=2.5cm,headheight=16pt]{geometry}
\usepackage{graphicx}
\usepackage{array}
\usepackage{booktabs}
\usepackage{tabularx}
\usepackage{multirow}
\usepackage{enumitem}
\usepackage{float}
\usepackage{titlesec}
\usepackage{fancyhdr}
\usepackage{setspace}
\usepackage{xcolor}
\usepackage{amsmath}

\usepackage{bidi}
\usepackage[no-math]{fontspec}
\usepackage[quiet]{polyglossia}

\setmainfont[Scale=1]{NotoKufiArabic-Regular.ttf}
\newfontfamily\arabicfont[Script=Arabic,Scale=0.95]{NotoKufiArabic-Regular.ttf}
\newfontfamily\englishfont[Scale=0.95]{TeX Gyre Termes}

\setdefaultlanguage[calendar=gregorian,hijricorrection=1]{arabic}
\setotherlanguage[variant=british]{english}

\usepackage[breaklinks,hidelinks]{hyperref}

\addto\captionsarabic{%
  \renewcommand{\contentsname}{الفهرس}%
  \renewcommand{\figurename}{الشكل}%
  \renewcommand{\tablename}{الجدول}%
}

\titleformat{\section}{\Large\bfseries}{\thesection}{0.7em}{}
\titlespacing*{\section}{0pt}{1.4ex plus 0.4ex}{0.9ex plus 0.3ex}

\pagestyle{fancy}
\fancyhf{}
\fancyhead[R]{\small\itshape سجل القرارات}
\fancyhead[L]{\small اسم المشروع}
\fancyfoot[C]{\small\thepage}
\renewcommand{\headrulewidth}{0.3pt}

\setlist[itemize]{topsep=0.3em, itemsep=0.15em, parsep=0pt}
\setlist[enumerate]{topsep=0.3em, itemsep=0.15em, parsep=0pt}

\newcolumntype{L}[1]{>{\raggedright\arraybackslash}p{#1}}
\newcolumntype{C}[1]{>{\centering\arraybackslash}p{#1}}

\setstretch{1.3}

\begin{document}

% ============================================================
% TITLE BLOCK
% ============================================================
\begin{center}
{\LARGE\bfseries سجل القرارات}\\[0.3em]
{\large Decision Log / Architecture Decision Records (ADR)}\\[1em]
\end{center}

% ============================================================
% PROJECT INFO
% ============================================================
% TODO: املأ بيانات المشروع.
\begin{table}[H]
\centering
\renewcommand{\arraystretch}{1.4}
\begin{tabularx}{\textwidth}{L{3.5cm} X L{3.5cm} X}
\toprule
\textbf{اسم المشروع:} & --- & \textbf{مدير المشروع:} & --- \\
\textbf{اسم الشركة:} & --- & \textbf{تاريخ التحديث:} & --- \\
\bottomrule
\end{tabularx}
\end{table}

\vspace{1em}

% ============================================================
% 1. DECISION SUMMARY TABLE
% ============================================================
\section{ملخص القرارات}
\label{sec:summary}

% TODO: أضف أو عدّل صفوف القرارات حسب مشروعك.
\begin{table}[H]
\centering
\caption{سجل القرارات}
\renewcommand{\arraystretch}{1.4}
\small
\begin{tabularx}{\textwidth}{C{1cm} L{1.8cm} L{3.5cm} L{3cm} L{2.5cm} C{1.5cm}}
\toprule
\textbf{المعرف} & \textbf{التاريخ} & \textbf{القرار} & \textbf{السياق} & \textbf{البديل المعتبر} & \textbf{الحالة} \\
\midrule
\LR{DEC-001} & --- & اعتماد المنهجية الرشيقة لإدارة المشروع & الحاجة إلى مرونة في التطوير والتسليم التدريجي & الشلال التقليدي & مقبول \\
\LR{DEC-002} & --- & اختيار قاعدة بيانات علائقية مفتوحة المصدر & متطلبات تخزين البيانات وتكامل الأنظمة & قواعد \LR{NoSQL} & مقبول \\
\LR{DEC-003} & --- & استخدام بنية الخدمات المصغّرة & قابلية التوسع وصيانة المكونات بشكل مستقل & البنية المتجانسة & قيد التنفيذ \\
\LR{DEC-004} & --- & تأجيل إطلاق الميزة الجديدة & عدم اكتمال الاختبارات ومخاطر الجودة & الإطلاق في الموعد المحدد & مقبول \\
\LR{DEC-005} & --- & تعيين مدير جودة مستقل & ضمان استقلالية مراجعة الجودة & دمج المسؤولية مع فريق التطوير & مقبول \\
\bottomrule
\end{tabularx}
\end{table}

% ============================================================
% 2. DECISION DETAILS
% ============================================================
\section{تفاصيل القرارات}
\label{sec:details}

% TODO: وسّع تفاصيل كل قرار حسب الحاجة.

\subsection{القرار \LR{DEC-001}: اعتماد المنهجية الرشيقة}
\noindent\textbf{التاريخ:} --- \quad \textbf{صاحب القرار:} --- \quad \textbf{الحالة:} مقبول \\[0.5em]
\textbf{السياق:} يحتاج المشروع إلى مرونة في التطوير مع التسليم التدريجي للمخرجات، مما يتطلب منهجية تتفاعل مع المتغيرات بسرعة. \\[0.3em]
\textbf{البدائل المعتبرة:} منهجية الشلال التقليدية، منهجية كانبان، منهجية سكرم. \\[0.3em]
\textbf{القرار:} اعتماد منهجية سكرم لإدارة المشروع مع دورات تسليم مدتها أسبوعان. \\[0.3em]
\textbf{النتائج:} تحسين سرعة الاستجابة للمتطلبات المتغيرة وزيادة شفافية التقدم.

\subsection{القرار \LR{DEC-002}: اختيار قاعدة بيانات علائقية}
\noindent\textbf{التاريخ:} --- \quad \textbf{صاحب القرار:} --- \quad \textbf{الحالة:} مقبول \\[0.5em]
\textbf{السياق:} يتطلب المشروع تخزين بيانات منظمة مع ضمان السلامة المرجعية والتكامل مع الأنظمة الحالية. \\[0.3em]
\textbf{البدائل المعتبرة:} قواعد بيانات \LR{NoSQL} (\LR{MongoDB})، قواعد بيانات تجارية (\LR{Oracle}). \\[0.3em]
\textbf{القرار:} استخدام \LR{PostgreSQL} كقاعدة بيانات رئيسية للمشروع. \\[0.3em]
\textbf{النتائج:} تقليل تكاليف الترخيص وضمان التوافق مع معايير \LR{SQL}.

\subsection{القرار \LR{DEC-003}: استخدام بنية الخدمات المصغّرة}
\noindent\textbf{التاريخ:} --- \quad \textbf{صاحب القرار:} --- \quad \textbf{الحالة:} قيد التنفيذ \\[0.5em]
\textbf{السياق:} النظام الحالي يعاني من صعوبة في التوسع وصيانة المكونات بشكل مستقل. \\[0.3em]
\textbf{البدائل المعتبرة:} البنية المتجانسة، البنية الموجهة بالخدمات. \\[0.3em]
\textbf{القرار:} الانتقال تدريجياً إلى بنية الخدمات المصغّرة بدءاً بالوحدات الحرجة. \\[0.3em]
\textbf{النتائج:} تحسين قابلية التوسع وتسهيل الصيانة، مع زيادة في التعقيد التشغيلي.

\subsection{القرار \LR{DEC-004}: تأجيل إطلاق الميزة الجديدة}
\noindent\textbf{التاريخ:} --- \quad \textbf{صاحب القرار:} --- \quad \textbf{الحالة:} مقبول \\[0.5em]
\textbf{السياق:} لم تكتمل اختبارات الجودة للميزة الجديدة، مما يزيد من مخاطر الإطلاق. \\[0.3em]
\textbf{البدائل المعتبرة:} الإطلاق في الموعد المحدد مع مخاطر معروفة، الإطلاق الجزئي. \\[0.3em]
\textbf{القرار:} تأجيل الإطلاق لمدة أسبوعين حتى اكتمال الاختبارات. \\[0.3em]
\textbf{النتائج:} تأخر الإطلاق مع ضمان جودة أعلى وتقليل مخاطر ما بعد الإطلاق.

\subsection{القرار \LR{DEC-005}: تعيين مدير جودة مستقل}
\noindent\textbf{التاريخ:} --- \quad \textbf{صاحب القرار:} --- \quad \textbf{الحالة:} مقبول \\[0.5em]
\textbf{السياق:} الحاجة إلى ضمان استقلالية مراجعة الجودة عن فريق التطوير. \\[0.3em]
\textbf{البدائل المعتبرة:} دمج مسؤولية الجودة مع فريق التطوير، الاستعانة بجهة خارجية. \\[0.3em]
\textbf{القرار:} تعيين مدير جودة مستقل داخل الفريق يرفع تقاريره مباشرة لإدارة المشروع. \\[0.3em]
\textbf{النتائج:} تحسين مستوى المراجعة وزيادة الثقة في جودة المخرجات.

\end{document}
